\documentclass[11pt]{article}
\usepackage[margin=1in]{geometry}
\usepackage{amsmath,amssymb,amsthm}
\usepackage{algorithm}
\usepackage{algpseudocode}
\usepackage{booktabs}
\usepackage{hyperref}

\newtheorem{definition}{Definition}
\newtheorem{theorem}{Theorem}
\newtheorem{lemma}[theorem]{Lemma}
\newtheorem{property}{Property}

\title{\textbf{Mathematical Specification, Spectral Acceleration, and Lookahead Dynamics of Heavy-Ball Momentum and Nesterov Accelerated Gradient (ALGO-NN-32)}}
\author{Team Scaibu \\ \small Email: \texttt{jaydeep.9664920749@gmail.com} \\ \small Architecture and Core Engine Team}
\date{\today}

\begin{document}
\maketitle

\begin{abstract}
First-order optimization on ill-conditioned anisotropic quadratics exhibits severe zig-zag oscillations and slow progress along low-curvature ravines. Polyak Heavy-Ball Momentum and Nesterov Accelerated Gradient (NAG) accumulate velocity vectors $\mathbf{v}_t$ to accelerate along consistent descent paths while canceling orthogonal oscillations. We formulate recurrence relations, prove the $\mathcal{O}(1/k^2)$ accelerated convergence bound, provide line-by-line pseudocode matching the reference implementation, and derive complexity.
\end{abstract}

\section{Notation and Mathematical Preliminaries}
\label{sec:notation}

Let $\theta \in \mathbb{R}^P$ denote parameter coordinates, $\mathbf{g} \in \mathbb{R}^P$ the gradient vector, and $\mathbf{v} \in \mathbb{R}^P$ the momentum velocity buffer. Let $\eta \in \mathbb{R}_{>0}$ denote learning rate and $\beta \in [0, 1)$ the momentum damping factor.

\subsection{Recurrence Formulations}
\paragraph{Intuitive Meaning:} Classical momentum acts like a heavy ball rolling down a hill, gaining speed in persistent directions. Nesterov momentum evaluates the gradient not at the current point, but at a projected lookahead position, applying a corrective braking force if the ball is moving uphill.

\paragraph{Formal Mathematical Definition:}
\begin{itemize}
    \item \textbf{Polyak Heavy-Ball Momentum:}
    \begin{equation}
    \label{eq:polyak_momentum}
    \mathbf{v}_{t+1} = \beta \mathbf{v}_t + \mathbf{g}_t, \quad \theta_{t+1} = \theta_t - \eta \mathbf{v}_{t+1}.
    \end{equation}
    \item \textbf{Nesterov Accelerated Gradient (NAG):}
    \begin{equation}
    \label{eq:nesterov_momentum}
    \mathbf{v}_{t+1} = \beta \mathbf{v}_t + \mathbf{g}_t, \quad \theta_{t+1} = \theta_t - \eta (\mathbf{g}_t + \beta \mathbf{v}_{t+1}).
    \end{equation}
\end{itemize}

\paragraph{Worked Numerical Example:}
Let $\theta = [1.0], \mathbf{g} = [0.2], \mathbf{v} = [0.1], \eta = 0.1, \beta = 0.9$:
$\mathbf{v}_{\mathrm{next}} = 0.9(0.1) + 0.2 = 0.09 + 0.2 = 0.29$.
Standard Momentum: $\Delta\theta = 0.1(0.29) = 0.029 \implies \theta_{\mathrm{next}} = 1.0 - 0.029 = 0.971$.
Nesterov: $\mathbf{u} = 0.2 + 0.9(0.29) = 0.2 + 0.261 = 0.461 \implies \Delta\theta = 0.1(0.461) = 0.0461 \implies \theta_{\mathrm{next}} = 1.0 - 0.0461 = 0.9539$.

\subsection{Summary of Symbols}
\begin{itemize}
    \item $P \in \mathbb{N}$: Parameter dimension.
    \item $\boldsymbol{\theta} \in \mathbb{R}^P$: Current parameter vector.
    \item $\mathbf{g} \in \mathbb{R}^P$: Gradient vector.
    \item $\mathbf{v} \in \mathbb{R}^P$: Momentum velocity vector.
    \item $\eta \in \mathbb{R}_{>0}$: Learning rate step size.
    \item $\beta \in [0, 1)$: Momentum decay coefficient.
    \item $\mathrm{nesterov} \in \{\text{True}, \text{False}\}$: Nesterov acceleration mode flag.
\end{itemize}

\section{Problem Definition and Core Concepts}
\label{sec:problem_definition}

\subsection{Spectral Radius and Condition Number Damping}
On an anisotropic quadratic $f(\theta) = \frac{1}{2} \theta^\top \mathbf{H} \theta$ with condition number $\kappa = \frac{\lambda_{\max}}{\lambda_{\min}}$, standard gradient descent converges at rate $\frac{\kappa - 1}{\kappa + 1}$, while momentum improves the convergence factor to $\frac{\sqrt{\kappa} - 1}{\sqrt{\kappa} + 1}$, achieving quadratic speedup.

\section{Algorithmic Specification}
\label{sec:algorithm}

Algorithm~\ref{alg:momentum_nag} specifies the parameter update procedure matching \texttt{impl.py}.

\begin{algorithm}[H]
\caption{Polyak Heavy-Ball and Nesterov Accelerated Gradient Update}
\label{alg:momentum_nag}
\begin{algorithmic}[1]
\Require Parameters $\boldsymbol{\theta} \in \mathbb{R}^P$, gradients $\mathbf{g} \in \mathbb{R}^P$, velocity $\mathbf{v} \in \mathbb{R}^P$, learning rate $\eta > 0$, momentum $\beta \in [0, 1)$, flag $\mathrm{nesterov}$.
\Ensure Updated parameters $\boldsymbol{\theta}_{\mathrm{next}} \in \mathbb{R}^P$, updated velocity $\mathbf{v}_{\mathrm{next}} \in \mathbb{R}^P$.
\State Initialize $\boldsymbol{\theta}_{\mathrm{next}} \gets [], \quad \mathbf{v}_{\mathrm{next}} \gets []$
\For{$i = 1$ to $P$}
    \State $v_{\mathrm{new}} \gets \beta \cdot v_i + g_i$
    \If{$\mathrm{nesterov} = \text{True}$}
        \State $step \gets \eta \cdot (g_i + \beta \cdot v_{\mathrm{new}})$
    \Else
        \State $step \gets \eta \cdot v_{\mathrm{new}}$
    \EndIf
    \State $\theta_{\mathrm{new}} \gets \theta_i - step$
    \State Append $\theta_{\mathrm{new}}$ to $\boldsymbol{\theta}_{\mathrm{next}}$, append $v_{\mathrm{new}}$ to $\mathbf{v}_{\mathrm{next}}$
\EndFor
\State \Return $\boldsymbol{\theta}_{\mathrm{next}}, \quad \mathbf{v}_{\mathrm{next}}$
\end{algorithmic}
\end{algorithm}

\section{Theoretical Guarantees and Acceleration Theorems}
\label{sec:guarantees}

\begin{theorem}[Nesterov Optimal Accelerated Convergence Rate]
\label{thm:nesterov_rate}
\textbf{Assumptions:} Let $f: \mathbb{R}^P \to \mathbb{R}$ be convex with $L$-Lipschitz continuous gradient $\nabla f$.
\textbf{Guarantee:} Nesterov acceleration achieves objective error $f(\theta_k) - f^* \le \frac{2 L \|\theta_0 - \theta^*\|_2^2}{(k+1)^2} = \mathcal{O}\left(\frac{1}{k^2}\right)$, matching the theoretical lower bound for first-order black-box optimization on smooth convex functions.
\textbf{Proof:}
Constructing the estimate sequence $\phi_k(\theta) = f(\theta_k) + \frac{\gamma_k}{2}\|\theta - v_k\|_2^2$ and applying induction over step $k$, the potential function decreases quadratically with momentum weight $\beta_k = \frac{k-1}{k+2}$. Standard gradient descent achieves $\mathcal{O}(1/k)$, establishing the strict superiority of Nesterov acceleration.
\textbf{Limitations:} In highly stochastic mini-batch settings, excessive momentum can amplify gradient variance, requiring careful warmup schedules.
\end{theorem}

\section{Complexity Analysis}
\label{sec:complexity}

\begin{itemize}
    \item \textbf{Time Complexity:} $\mathcal{O}(P)$ sequential floating-point multiply-accumulate operations.
    \item \textbf{Space Complexity:} $\mathcal{O}(P)$ memory for the velocity buffer $\mathbf{v}$ (1 extra tensor of parameter size).
\end{itemize}

\section{Worked Numerical Example}
\label{sec:worked_example}

Let $\theta = [0.0], \mathbf{g} = [1.0], \mathbf{v} = [0.0], \eta = 0.01, \beta = 0.9$ across two steps:
\begin{enumerate}
    \item Step 1: $v_1 = 0.9(0) + 1.0 = 1.0 \implies \theta_1 = 0.0 - 0.01(1.0) = -0.01$.
    \item Step 2 (assuming $\mathbf{g} = 1.0$ again): $v_2 = 0.9(1.0) + 1.0 = 1.9 \implies \theta_2 = -0.01 - 0.01(1.9) = -0.029$.
    Notice the effective step size doubled from $0.01$ to $0.019$ due to momentum velocity build-up.
\end{enumerate}

\section{Numerical Considerations and Edge Cases}
\label{sec:numerical}

\begin{itemize}
    \item \textbf{Effective Learning Rate Rescaling:} When tuning $\beta$, the steady-state effective learning rate is $\eta_{\mathrm{eff}} = \frac{\eta}{1 - \beta}$; scaling $\beta$ from $0.9$ to $0.99$ increases effective step magnitude by $10\times$.
\end{itemize}

\begin{thebibliography}{9}
\bibitem{nesterov1983} Y.~Nesterov, ``A Method for Solving the Convex Programming Problem with Convergence Rate $O(1/k^2)$,'' \emph{Doklady AN SSSR}, vol.~269, pp.~543--547, 1983.
\bibitem{polyak1964} B.~T.~Polyak, ``Some Methods of Speeding Up the Convergence of Iteration Methods,'' \emph{USSR Comput. Math. Math. Phys.}, vol.~4, no.~5, pp.~1--17, 1964.
\end{thebibliography}

\end{document}
